%! TEX root = ../main.tex

\section{Лекция от 28.03.2023}
Нейрон должен обладать входом и выходом, а также силой связи с соседними
нейронами. Также должна быть функция активации, и она должна быть простой.
Общая формула искусственного нейрона:
\begin{center}
  \begin{tikzpicture}[node distance = 2cm]
    \node[circle, draw = black] (s) {$S$};
    \node[circle, draw=black, minimum size=0.5cm, right of = s] (out) {};
  \end{tikzpicture}
\end{center}
\[
  S = \sum_{i=1}^{n} X_i \omega_i + \omega_0
\] 
\[
  step(x) = \left\{\begin{aligned} 
    &1, x > 0 \\
    &0, x \le 0
  \end{aligned}\right. 
\] 
\begin{center}
  \begin{tikzpicture}[node distance = 3cm]
    \node[block] (n) {$+\omega_0$}; 
    \node[left of = n] (l) {$x$};
    \node[right of = n] (r) {Выход};
    \draw[arrow] (l) -- node[anchor=south]{$ \omega $} (n);
    \draw[arrow] (n) -- node[anchor=south]{$f$} (r);
  \end{tikzpicture}
\end{center}

Сигмойда:
\[
  \sigma = \frac{e^{x}}{1 + e^{x}}
\] 

Гиперболический тангенс:
\[
  th(x) = \frac{sh(x)}{ch(x)} = \frac{e^{x} - e^{-x}}{e^{x} + e^{-x}}
\] 

ReLU:
\[
  a  
\] 

Пример искусственного нейрона:
\begin{center}
  \begin{tikzpicture}[node distance = 3cm]
    \node[block] (n) {$ \omega _0$}; 
    \node[rectangle, draw, left of = n] (x1) {$x_1$};
    \node[right of = n] (out) {Выход};
    \node[rectangle, draw, below of = x1, yshift=2cm] (x2) {$x_2$};
    \draw[arrow] (x1) -- node[anchor=south]{$\omega$}(n);
    \draw[arrow] (n) -- node[anchor=south]{$f$}(out);
    \draw[arrow] (x2) -- (n);
  \end{tikzpicture}
\end{center}

{\Huge\textbf{Нейронные сети}}

Задача:\par


План решения регрессии для нейросети:
\begin{enumerate}
  \item Взять тренировочную выборку
\end{enumerate}
